\chapter{Реализация}

Все звенья контура реализованы и разворачиваются с нуля из репозитория одной
командой. Объём собственного кода и конфигурации составляет около 14~тысяч
строк без учёта тестов.

\section{Экспериментальный стенд}

Стенд состоит из трёх кластеров kind~\cite{kind} на одной машине (Apple M4,
16~ГБ оперативной памяти, из них 10~ГБ выделено виртуальной машине
контейнеров): управляющего кластера с компонентами контура и двух кластеров
федерации. Третий кластер федерации не поместился в доступную память; все
сценарии экспериментов выполнимы на двух. Кластеры находятся в общей
docker-сети и обращаются друг к другу по внутренним адресам.

Тестовый сервис реализован на Go и состоит из двух звеньев: frontend
вызывает backend с таймаутом 1~с и при ошибке или таймауте отвечает кодом
502. Благодаря этому отказ, внесённый в backend, отражается в метриках
frontend, по которым каталог оценивает кластер. В каждом кластере федерации
работает долгоживущий сервис \texttt{app}, являющийся объектом переключения
трафика, и поток короткоживущих нагрузок (workload): новая нагрузка
появляется каждые 30~с и живёт 180~с, так что одновременно существует около
шести нагрузок. Весь трафик проходит через шлюз Envoy
Gateway~\cite{envoy-gateway} в управляющем кластере; веса маршрутов задаются
через Gateway API~\cite{gateway-api}.

\section{Наблюдаемость и трассировка}

В каждом кластере федерации работают Prometheus (опрос целей раз в 5~с) и
Thanos sidecar. В управляющем кластере Thanos Querier служит единой точкой
запросов ко всем кластерам, Thanos Ruler вычисляет глобальные правила и
алерты, Grafana визуализирует данные. Объектное хранилище не используется:
долгосрочным хранилищем результатов служит выгрузка сырых рядов каждого
прогона в репозиторий.

Локальные показатели (SLI) кластера~--- доля ошибок и 95-й процентиль
задержки frontend за 5~минут (для каталога) и за 1~минуту (для контроллера
восстановления), а также загрузка узла как максимум загрузки процессора и
памяти. Глобальный показатель доступности федерации вычисляется на шлюзе,
а не как сумма серверных показателей кластеров; причины этого решения
обсуждаются в главе~\ref{ch:discussion}.

Трассировка построена на OpenTelemetry~\cite{otel} и Jaeger~\cite{jaeger}.
Корнем трейса является шлюз, он же выбирает 10\% запросов для записи и
передаёт контекст трассировки~\cite{trace-context} в кластер. Это позволяет
разложить задержку запроса на звенья: шлюз, переход из управляющего кластера
в кластер федерации, frontend, сеть до backend и backend.

\section{Каталог кластеров}

Каталог~--- REST-сервис на FastAPI~\cite{fastapi} и
PostgreSQL~\cite{postgresql}, работающий в управляющем кластере. Раз в 30~с
каталог запрашивает метрики каждого кластера у Thanos, записывает снимок и
пересчитывает статус кластера: \texttt{degraded} при доле ошибок выше 5\% или
задержке p95 выше 200~мс, \texttt{unknown} при отсутствии метрик. Раз в 10~с
фоновый процесс читает события и отчёты Kyverno во всех кластерах и
записывает обнаруженные нарушения политик с указанием допустимых кластеров.
Компоненты аутентифицируются статическими токенами; при регистрации
кластера каталог сверяет переданные атрибуты с описанием, хранящимся в
самом кластере.

\section{Политики размещения}

Политики исполняются на двух независимых рубежах. Первый рубеж~--- фильтрация
кандидатов каталогом по региону, сертификациям и ресурсам. Второй рубеж~---
две политики Kyverno в каждом кластере, проверяющие соответствие требований
нагрузки к региону и сертификациям фактическим атрибутам кластера.
Ресурсные ограничения внутри кластера обеспечивает стандартный планировщик
Kubernetes. Первый рубеж является оптимизацией, второй~--- гарантией: он
перехватывает ошибки планировщика, размещения в обход него и рассогласование
каталога с реальным состоянием кластеров.

\section{Планировщик уровня федерации}
\label{sec:scheduler}

Планировщик реализован на Go с использованием библиотеки
controller-runtime~\cite{controller-runtime}. Входом служит объект
\texttt{Workload} в управляющем кластере. Планировщик получает из каталога
список допустимых кластеров со снимками их метрик, выбирает кластер,
записывает решение со всеми кандидатами, создаёт нагрузку в выбранном кластере
и маршрут на шлюзе. Отказ политики при создании нагрузки фиксируется в
каталоге как нарушение.

\subsection{Оценка кандидата}

Каждому допустимому кластеру ставится оценка
\begin{equation}
  \label{eq:score}
  \mathrm{score} = w_1 (1 - e) + w_2 \left(1 - \min\left(\frac{p_{95}}{L}, 1\right)\right)
                 + w_3 (1 - s),
\end{equation}
где $e$~--- доля ошибок, $p_{95}$~--- 95-й процентиль задержки, $s$~---
загрузка узла, $L = 200$~мс~--- целевое значение задержки,
$w_1 = 0.5$, $w_2 = 0.3$, $w_3 = 0.2$~--- веса. Каждое слагаемое принимает
значения от 0 (плохо) до 1 (хорошо), сумма весов равна единице, поэтому и
оценка лежит в отрезке $[0, 1]$. Задержка нормируется на целевое значение:
любая задержка выше 200~мс считается одинаково неприемлемой. Наибольший вес
получает доля ошибок как признак, непосредственно наблюдаемый пользователем.
Веса зафиксированы до проведения экспериментов, и целевой показатель
вычисляется только с ними.

Выбирается кластер с наибольшей оценкой; при равенстве оценок (разница
меньше $10^{-6}$) кластеры выбираются по очереди. Отсутствующая или
устаревшая (старше 60~с) метрика даёт своему слагаемому нейтральное значение
0.5. Если метрик нет ни у одного кандидата, все оценки равны, и планировщик
в точности совпадает с наивным: преимущество возникает только при наличии
сигнала.

\subsection{Пример работы}

В табл.~\ref{tab:scoring-example} приведены решения планировщика в одном из
прогонов сценария деградации сети. Отказ (задержка 500~мс и потеря 10\%
пакетов до backend) внесён в кластер B в момент 0 и снят через 300~с; новая
нагрузка поступает каждые 30~с.

\begin{table}[ht]
  \centering
  \caption{Оценки кластеров при деградации сети в кластере B}
  \label{tab:scoring-example}
  \small
  \begin{tabular}{@{}rlccc@{}}
    \toprule
    Время, с & Состояние кластера B & Оценка A & Оценка B & Выбор \\
    \midrule
    $-64$ & норма & 0.939 & 0.935 & A \\
    $-34$ & норма & 0.930 & 0.936 & B \\
    $-4$  & норма & 0.929 & 0.939 & B \\
    $26$  & p95 679 мс, ошибок 1.5\% & 0.941 & 0.635 & A \\
    $116$ & p95 770 мс, ошибок 4.7\% & 0.938 & 0.621 & A \\
    $266$ & p95 1660 мс, ошибок 8.8\% & 0.931 & 0.603 & A \\
    $566$ & отказ снят, окно 5 мин ещё не очищено & 0.935 & 0.619 & A \\
    $596$ & норма & 0.931 & 0.937 & B \\
    \bottomrule
  \end{tabular}
\end{table}

В штатном режиме оценки кластеров почти равны, и выбор определяется
незначительными колебаниями загрузки. Через 26~с после внесения отказа
оценка кластера B падает с 0.94 до 0.63. Падение почти полностью объясняется
задержкой: p95 превысил 200~мс, и второе слагаемое обнулилось, уменьшив
оценку на 0.3, тогда как доля ошибок в этот момент составляла лишь 1.5\% и
уменьшила оценку на 0.008. Таким образом, при сетевой деградации основным
сигналом служит задержка. Разрыв оценок (около 0.3) на порядок превышает их
колебания в штатном режиме (около 0.01), поэтому выбор однозначен. После
снятия отказа кластер B ещё около пяти минут сохраняет низкую оценку, пока
неблагоприятные значения не покинут пятиминутное окно; это обеспечивает
естественную осторожность при возврате нагрузки.

\subsection{Наивный планировщик}

Наивный планировщик выбирает кластеры по очереди (round-robin) среди тех же
допустимых кандидатов, полученных из каталога. Режим переключается
параметром запуска без пересборки, так что оба варианта~--- один и тот же
исполняемый файл. Наивный планировщик также записывает снимок метрик всех
кандидатов, что позволяет офлайн вычислить, какой выбор сделал бы
разрабатываемый планировщик при любых весах.

\section{Контроллер аварийного восстановления}

Контроллер реализован как оператор Kubernetes на Go. Раз в 10~с он снимает
по каждому кластеру три сигнала: пробу сервиса \texttt{app} по тому же
адресу, куда направляется трафик; пробу API-сервера кластера; долю ошибок за
1~минуту из Thanos. Состояние кластера ведёт конечный автомат
(рис.~\ref{fig:fsm}) с асимметричными переходами: деградация признаётся за
один цикл, а трафик снимается только после трёх неудачных проб сервиса
подряд и возвращается ступенями после серии успешных циклов. Асимметрия
защищает от многократного переключения трафика туда и обратно.

\begin{figure}[ht]
  \centering
  \begin{tikzpicture}[
      font=\scriptsize, >={Stealth[length=2mm]}, node distance=14mm and 30mm,
      st/.style={draw, rounded corners=3pt, minimum width=24mm, minimum height=8mm, font=\small\ttfamily},
      lbl/.style={align=center, inner sep=2pt}]
    \node[st] (h) {HEALTHY};
    \node[st, right=of h] (d) {DEGRADED};
    \node[st, below=22mm of d] (dn) {DOWN};
    \node[st, below=22mm of h] (r) {RECOVERING};
    \draw[->] ([yshift=2mm]h.east) -- ([yshift=2mm]d.west)
      node[lbl, midway, above] {любой сигнал\\не в норме, 1 цикл};
    \draw[->] ([yshift=-2mm]d.west) -- ([yshift=-2mm]h.east)
      node[lbl, midway, below] {3 чистых цикла};
    \draw[->] (d) -- (dn) node[lbl, midway, right] {проба сервиса\\неудачна\\3 цикла подряд};
    \draw[->] ([yshift=2mm]dn.west) -- ([yshift=2mm]r.east)
      node[lbl, midway, above] {сервис отвечает +\\контрольная серия};
    \draw[->] ([yshift=-2mm]r.east) -- ([yshift=-2mm]dn.west)
      node[lbl, midway, below] {проба неудачна\\3 цикла подряд};
    \draw[->] (r) -- (h) node[lbl, midway, left] {вес 10\% $\to$ 50\% $\to$ 100\%,\\по 2 чистых цикла\\на ступень};
  \end{tikzpicture}
  \caption{Конечный автомат состояния кластера в контроллере восстановления}
  \label{fig:fsm}
\end{figure}

При переходе кластера в состояние DOWN контроллер обнуляет его вес на
маршруте шлюза и устанавливает статус \texttt{down} в каталоге. При
восстановлении контроллер выполняет контрольную серию из 10 запросов и
возвращает вес ступенями 10\%, 50\% и 100\%. Действия выводятся из
текущего состояния на каждом цикле, поэтому неудачная запись или перезапуск
контроллера исправляются на следующем цикле. Если в состоянии DOWN находятся
все кластеры, веса остаются номинальными: обнулять последний доступный
маршрут бессмысленно. Режим наблюдения ведёт тот же автомат и записывает
инциденты, но трафик не переключает; он служит базовой линией «без
аварийного восстановления».
